% GENERATED FILE. Edit canonical lessons, then run npm run generate:books.
% canonical-source-hash: fnv1a64:8f116c173ca31b36
% canonical-lessons: RU-C165-vorona, RU-C165-oryol, RU-C165-popugay, RU-C165-khomyak, RU-C165-yozh

\chapter{Crow, Eagle, Parrot, Hamster, Hedgehog}
\label{ch:ru-crow-eagle-parrot-hamster-hedgehog}
\hlchaptermodality{165}

\begin{chapteropening}
\textbf{By the end of this chapter:} \emph{I can say a crow, an eagle, a parrot, a hamster and a hedgehog.}

Complete the last lesson of chapter 165: I can say a crow, an eagle, a parrot, a hamster and a hedgehog.
\end{chapteropening}

\section[\texorpdfstring{\ru{ворона} (voróna)}{ворона (voróna)}]{\ru{ворона} (voróna) — a crow}

\label{lesson:RU-C165-vorona}



\begin{quote}
Before the new one: say the Russian for air, then the Russian for the earth; land.
\end{quote}

\subsection*{You'll want to know: \ru{ворона}}
\textbf{\ru{ворона}} (\emph{voróna}) — \textquotedblleft{}a crow\textquotedblright{}.

\begin{cousinweb}
\textbf{\ru{ворона}} $\leftarrow$ Proto-Slavic \emph{*vorna} — inherited, a cousin of Lithuanian \emph{várna} and Old Prussian \emph{warne}, \textquotedblleft{}crow\textquotedblright{} (Vasmer); it goes with \emph{\ru{ворон}}, \textquotedblleft{}raven\textquotedblright{}. Church Slavonic has \emph{\ru{врана}}, \textquotedblleft{}crow\textquotedblright{}.
\end{cousinweb}

\begin{grammarlens}[title={What you've built}]
One more word for this chapter.
\end{grammarlens}

\subsection*{Your turn}
\begin{itemize}
\raggedright
  \item \emph{Say it:} \emph{voróna}
  \item \emph{Say it:} \emph{voróna}, once more
  \item \emph{Say it:} say \emph{voróna}
  \item \emph{Recall:} say the Russian for air, then the Russian for the earth; land, then say \emph{voróna} again
\end{itemize}

\begin{tcolorbox}[breakable,colback=teal!4,colframe=teal!35!black,title={Before you move on}]
What does \ru{ворона} mean? (\textquotedblleft{}A crow\textquotedblright{}.) Say it once more.
\end{tcolorbox}
\section[\texorpdfstring{\ru{орёл} (oryól)}{орёл (oryól)}]{\ru{орёл} (oryól) — an eagle}

\label{lesson:RU-C165-oryol}



\begin{quote}
Before the new one: say the Russian for the earth; land, then the Russian for a crow.
\end{quote}

\subsection*{You'll want to know: \ru{орёл}}
\textbf{\ru{орёл}} (\emph{oryól}) — \textquotedblleft{}an eagle\textquotedblright{}.

\begin{cousinweb}
\textbf{\ru{орёл}} $\leftarrow$ Proto-Slavic \emph{*orĭlŭ} — inherited. Vasmer lists Lithuanian \emph{erẽlis}, Gothic \emph{ara} and German \emph{Aar}, \textquotedblleft{}eagle\textquotedblright{}, and Greek \emph{órnis}, \textquotedblleft{}bird\textquotedblright{}; English \emph{erne}, \textquotedblleft{}sea eagle\textquotedblright{}, is a cousin.
\end{cousinweb}

\begin{grammarlens}[title={What you've built}]
One more word for this chapter.
\end{grammarlens}

\subsection*{Your turn}
\begin{itemize}
\raggedright
  \item \emph{Say it:} \emph{oryól}
  \item \emph{Say it:} \emph{oryól}, once more
  \item \emph{Say it:} say \emph{oryól}
  \item \emph{Recall:} say the Russian for the earth; land, then the Russian for a crow, then say \emph{oryól} again
\end{itemize}

\begin{tcolorbox}[breakable,colback=teal!4,colframe=teal!35!black,title={Before you move on}]
What does \ru{орёл} mean? (\textquotedblleft{}An eagle\textquotedblright{}.) Say it once more.
\end{tcolorbox}
\section[\texorpdfstring{\ru{попугай} (popugáy)}{попугай (popugáy)}]{\ru{попугай} (popugáy) — a parrot}

\label{lesson:RU-C165-popugay}



\begin{quote}
Before the new one: say the Russian for a crow, then the Russian for an eagle.
\end{quote}

\subsection*{You'll want to know: \ru{попугай}}
\textbf{\ru{попугай}} (\emph{popugáy}) — \textquotedblleft{}a parrot\textquotedblright{}.

\begin{cousinweb}
\textbf{\ru{попугай}} — recorded from 1589; it probably came through Dutch \emph{papegaai} and Low German \emph{papagoie} from Old French \emph{papegai} and Spanish \emph{papagayo}, from Arabic \emph{babaghā}, \textquotedblleft{}parrot\textquotedblright{} (Vasmer).
\end{cousinweb}

\begin{grammarlens}[title={What you've built}]
One more word for this chapter.
\end{grammarlens}

\subsection*{Your turn}
\begin{itemize}
\raggedright
  \item \emph{Say it:} \emph{popugáy}
  \item \emph{Say it:} \emph{popugáy}, once more
  \item \emph{Say it:} say \emph{popugáy}
  \item \emph{Recall:} say the Russian for a crow, then the Russian for an eagle, then say \emph{popugáy} again
\end{itemize}

\begin{tcolorbox}[breakable,colback=teal!4,colframe=teal!35!black,title={Before you move on}]
What does \ru{попугай} mean? (\textquotedblleft{}A parrot\textquotedblright{}.) Say it once more.
\end{tcolorbox}
\section[\texorpdfstring{\ru{хомяк} (khomyák)}{хомяк (khomyák)}]{\ru{хомяк} (khomyák) — a hamster}

\label{lesson:RU-C165-khomyak}



\begin{quote}
Before the new one: say the Russian for an eagle, then the Russian for a parrot.
\end{quote}

\subsection*{You'll want to know: \ru{хомяк}}
\textbf{\ru{хомяк}} (\emph{khomyák}) — \textquotedblleft{}a hamster\textquotedblright{}.

\begin{cousinweb}
\textbf{\ru{хомяк}} — Old Russian \emph{\ru{хомекы}}; Vasmer thinks it is probably an old loan from Iranian, compare Avestan \emph{hamaēstar-}, \textquotedblleft{}a foe who throws to the ground\textquotedblright{}.
\end{cousinweb}

\begin{grammarlens}[title={What you've built}]
One more word for this chapter.
\end{grammarlens}

\subsection*{Your turn}
\begin{itemize}
\raggedright
  \item \emph{Say it:} \emph{khomyák}
  \item \emph{Say it:} \emph{khomyák}, once more
  \item \emph{Say it:} say \emph{khomyák}
  \item \emph{Recall:} say the Russian for an eagle, then the Russian for a parrot, then say \emph{khomyák} again
\end{itemize}

\begin{tcolorbox}[breakable,colback=teal!4,colframe=teal!35!black,title={Before you move on}]
What does \ru{хомяк} mean? (\textquotedblleft{}A hamster\textquotedblright{}.) Say it once more.
\end{tcolorbox}
\section[\texorpdfstring{\ru{ёж} (yozh)}{ёж (yozh)}]{\ru{ёж} (yozh) — a hedgehog}

\label{lesson:RU-C165-yozh}



\begin{quote}
Before the new one: say the Russian for a parrot, then the Russian for a hamster.
\end{quote}

\subsection*{You'll want to know: \ru{ёж}}
\textbf{\ru{ёж}} (\emph{yozh}) — \textquotedblleft{}a hedgehog\textquotedblright{}.

\begin{cousinweb}
\textbf{\ru{ёж}} $\leftarrow$ Proto-Slavic \emph{*ežĭ} — inherited. Vasmer lists Lithuanian \emph{ežỹs}, German \emph{Igel} and Greek \emph{echînos}, \textquotedblleft{}hedgehog\textquotedblright{}, and ties them to Greek \emph{échis}, \textquotedblleft{}snake\textquotedblright{}: perhaps a taboo name, \textquotedblleft{}the snake-eater\textquotedblright{}.
\end{cousinweb}

\begin{grammarlens}[title={What you've built}]
One more word for this chapter.
\end{grammarlens}

\subsection*{Your turn}
\begin{itemize}
\raggedright
  \item \emph{Say it:} \emph{yozh}
  \item \emph{Say it:} \emph{yozh}, once more
  \item \emph{Say it:} say \emph{yozh}
  \item \emph{Recall:} say the Russian for a parrot, then the Russian for a hamster, then say \emph{yozh} again
\end{itemize}

\begin{tcolorbox}[breakable,colback=teal!4,colframe=teal!35!black,title={Before you move on}]
What does \ru{ёж} mean? (\textquotedblleft{}A hedgehog\textquotedblright{}.) Say it once more.
\end{tcolorbox}
